\chapter{Production DataOps and the Enterprise Capstone}
\index{Terraform}\index{Bicep}\index{dbt}\index{CI/CD}\index{GitHub Actions}\index{expand--contract}\index{infrastructure as code}%
\begin{objectives}
\begin{itemize}
    \item Provision ADLS Gen2, Synapse, and Data Factory with Terraform's \texttt{azurerm} provider, and read Bicep as the Azure-native alternative.
    \item Manage Terraform state in an Azure Storage backend with locking, and keep state files out of version control.
    \item Structure a dbt project---models, tests, documentation---against Synapse or Fabric.
    \item Build a GitHub Actions pipeline that validates Terraform and runs dbt tests on every pull request, with \texttt{apply} gated on review.
    \item Deploy zero-downtime schema changes with the expand--contract pattern coordinated through dbt and lineage.
    \item Design environment promotion (dev $\rightarrow$ test $\rightarrow$ prod) for data platform artifacts.
    \item Complete the program's final capstone: a secure, governed, multi-subscription data mesh deployed entirely from code.
\end{itemize}
\end{objectives}

\begin{crosscloud}
Infrastructure as code, SQL-based transformation, and CI/CD are the most portable skills in this book---the same Terraform workflow and the same dbt project structure run against every cloud.

\vspace{2mm}
{\footnotesize
\begin{tabular}{@{}p{3.0cm}p{2.9cm}p{3.1cm}p{3.2cm}@{}}
\toprule
\textbf{Capability} & \textbf{AWS} & \textbf{Azure} & \textbf{GCP} \\
\midrule
Multi-cloud IaC & Terraform \texttt{aws} provider & Terraform \texttt{azurerm} provider & Terraform \texttt{google} provider \\
Native IaC & CloudFormation / CDK & Bicep / ARM templates & Deployment Manager (legacy) / Config Controller \\
Transform-as-code & dbt on Redshift/Athena & dbt on Synapse/Fabric & dbt on BigQuery \\
CI/CD & CodePipeline / GitHub Actions & GitHub Actions / Azure DevOps & Cloud Build / GitHub Actions \\
State backend & S3 + DynamoDB lock & Azure Storage + lease lock & GCS backend \\
\addlinespace
\multicolumn{4}{@{}l@{}}{\textbf{Fabric}: Git integration + deployment pipelines; \textbf{Snowflake}: dbt + Terraform provider; \textbf{MongoDB}: Atlas Terraform provider.} \\
\bottomrule
\end{tabular}}

\vspace{1mm}
Everything in this chapter is vendor-neutral discipline: declarative infrastructure, reviewed plans, tested transformations, and schema changes that never break a consumer mid-flight.
\end{crosscloud}

\begin{keyterms}
\textbf{IaC}: infrastructure as code---declarative, versioned, reviewable environment definitions.\quad
\textbf{Terraform state}: the recorded mapping of configuration to real resources; shared state needs a remote backend with locking.\quad
\textbf{dbt}: SQL transformation framework with models, tests, and documentation as code.\quad
\textbf{Expand--contract}: zero-downtime schema evolution: add new shape, migrate, switch reads, then remove the old.\quad
\textbf{Environment promotion}: the same artifacts progressing dev $\rightarrow$ test $\rightarrow$ prod through gates.
\end{keyterms}

\section{Week 24 Overview: The Program Becomes a Platform}
The final week is about \emph{repeatability}. Every resource built by hand in Chapters 1--23 becomes code; every deployment becomes a reviewed, tested pipeline run; every schema change becomes a migration that never breaks a consumer. This is DataOps: the application of software-engineering delivery discipline to data platforms. The chapter ends with the program's final capstone---the multi-subscription, zero-trust, governed data mesh that assembles everything.

\section{Monday: Terraform and Bicep for Data Platforms}
\index{Terraform}\index{Bicep}\index{state backend}

\subsection{Terraform with the azurerm Provider}
Terraform declares desired state; the \texttt{azurerm} provider maps it to Azure Resource Manager \cite{hashicorp2025azurerm}. The data-platform skeleton:

\begin{lstlisting}[language={},caption={Terraform for the storage and orchestration foundation.},label={lst:ch24-terraform}]
terraform {
  required_providers {
    azurerm = { source = "hashicorp/azurerm", version = "~> 4.0" }
  }
  backend "azurerm" {
    resource_group_name  = "rg-tfstate"
    storage_account_name = "sttfstate001"
    container_name       = "tfstate"
    key                  = "platform.tfstate"
  }
}
provider "azurerm" { features {} }

resource "azurerm_storage_account" "adls" {
  name                     = "stadlsdemo001"
  resource_group_name      = azurerm_resource_group.rg.name
  location                 = azurerm_resource_group.rg.location
  account_tier             = "Standard"
  account_replication_type = "LRS"
  is_hns_enabled           = true
  public_network_access_enabled = false
}

resource "azurerm_data_factory" "adf" {
  name                = "adf-demo-001"
  resource_group_name = azurerm_resource_group.rg.name
  location            = azurerm_resource_group.rg.location
  identity { type = "SystemAssigned" }
}
\end{lstlisting}

Two details carry the production weight. First, the \textbf{remote state backend} on Azure Storage with lease-based locking: state is the source of truth about what exists, so it must be shared, locked against concurrent applies, and never committed to git. Second, explicit security defaults: \texttt{public\_network\_access\_enabled = false} written \emph{in code}, because several resource types default to public.

\subsection{Bicep: The Azure-Native Alternative}
Bicep transpiles to ARM templates and offers first-day support for every Azure resource property, tight portal integration, and no state file (ARM tracks deployments) \cite{microsoft2025bicep}. Choose Bicep for Azure-only estates with ARM-native teams; choose Terraform for multi-cloud posture, the provider ecosystem (Databricks, dbt Cloud, Snowflake), and the plan/apply workflow maturity. Many organizations standardize on Terraform for data platforms and Bicep for foundational landing zones.

\begin{pitfall}
\textbf{terraform apply without review.} A destructive apply on a storage account is a data-loss event. The workflow is: \texttt{plan} on every pull request, review the plan like code, \texttt{apply} only from the protected main branch after merge---and state locking makes two applies at once impossible.
\end{pitfall}

\section{Tuesday: dbt---Transformations as Reviewed Code}
\index{dbt}\index{incremental models}

\subsection{Models, Tests, Documentation}
A dbt project is a directory of SQL \textbf{models} (SELECT statements materialized as views/tables/incremental tables), \textbf{tests} (generic---unique, not-null, accepted-values, relationships---plus custom SQL tests), and generated \textbf{documentation} with lineage \cite{dbtlabs2025core}. Against Synapse or Fabric, dbt becomes the transformation layer the warehouse was missing: versioned SQL, reviewed in PRs, tested in CI, documented automatically. The \texttt{ref()} graph is a dependency DAG you did not have to draw---and, as Wednesday shows, the safety net for schema changes.

\subsection{Incremental Models}
Full refreshes do not scale; dbt \textbf{incremental} models append or merge new rows using a watermark column. The pattern pairs with Chapter 12's idempotency: incremental models with unique keys merge rather than duplicate, and \texttt{--full-refresh} remains the recovery path when logic changes invalidate history.

\section{Wednesday: CI/CD and Zero-Downtime Schema Evolution}
\index{CI/CD}\index{GitHub Actions}\index{expand--contract}

\subsection{The Pipeline}
The data platform's CI/CD workflow on GitHub Actions: on pull request---\texttt{terraform fmt/validate/plan}, \texttt{dbt build} against a CI schema, linting; on merge to main---\texttt{terraform apply}, ADF publish (the npm ADFUtilities pattern), dbt deploy to the next environment. Environment promotion is the same artifact set advancing through dev $\rightarrow$ test $\rightarrow$ prod with gates between: plan review, test evidence, and for prod, an approval.

\begin{lstlisting}[language={},caption={Pull-request validation workflow (excerpt).},label={lst:ch24-actions}]
name: pr-validate
on: [pull_request]
jobs:
  terraform:
    runs-on: ubuntu-latest
    steps:
      - uses: actions/checkout@v4
      - uses: hashicorp/setup-terraform@v3
      - run: terraform -chdir=infra init -backend=false
      - run: terraform -chdir=infra validate
      - run: terraform -chdir=infra plan -input=false
  dbt:
    runs-on: ubuntu-latest
    steps:
      - uses: actions/checkout@v4
      - run: pip install dbt-synapse
      - run: dbt build --target ci        # models + tests against CI schema
\end{lstlisting}

\subsection{Expand--Contract Migrations}
A breaking schema change---renaming \texttt{customer\_email}---ships as four ordered releases: (1) \textbf{expand}: add the new nullable column alongside the old; (2) \textbf{migrate}: dual-write new writes and backfill history (dbt incremental or ADF copy); (3) \textbf{switch reads}: a backward-compatible view shim so consumers keep one stable interface; (4) \textbf{contract}: drop the old column only after Purview lineage and the dbt \texttt{ref()} graph prove no consumer reads it. Every intermediate state is valid; no consumer breaks between releases.

\begin{pitfall}
\textbf{Contracting first.} The failure story is canonical: a release dropped \texttt{customer\_email} while the dbt model and the Power BI dataset still referenced it; dashboards failed at 09:00 and the fix needed a hot rollback and a full dbt re-run. Never drop until the dependency check---lineage plus \texttt{ref()} graph---returns zero readers, and ship the view shim in the same release as the column change.
\end{pitfall}

\section{Thursday Hands-on Project: Terraform + dbt + GitHub Actions}
\index{hands-on lab}\index{Terraform!project}
Create a GitHub repository containing Terraform that provisions an ADLS Gen2 account and an ADF instance, a dbt project with three models and five tests, and a workflow that validates Terraform and runs dbt tests on every pull request.

\subsection{Build the Repository}
\begin{enumerate}
    \item \texttt{infra/}: the Terraform from Listing~\ref{lst:ch24-terraform}---resource group, HNS-enabled ADLS account with public access disabled, ADF with a system-assigned identity, and the remote state backend.
    \item \texttt{transform/}: a dbt project with staging, intermediate, and mart models over the lab orders data; five tests: unique and not-null on the order key, accepted values on status, a relationship test to customers, a freshness test on source arrival.
    \item \texttt{.github/workflows/}: Listing~\ref{lst:ch24-actions}, extended with the dbt job against a CI schema.
\end{enumerate}

\subsection{Prove the Gates}
Open three pull requests: a valid change (both jobs green, plan output posted), a Terraform change that would enable public access (plan review catches it---this is why humans read plans), and a dbt change violating a test (dbt job red, merge blocked by branch protection). The red PRs are the deliverables: gates that have never fired are gates you have not verified.

\section{Friday: Program Review and Certification Scheduling}
\index{DP-700!revision}
Close the program deliberately. Re-run Chapter 23's exam-mapping table against yourself; the capstones you shipped are the scenario practice the exam assumes. Identify the two weakest domains by mock-exam evidence, schedule both full-length mocks for the coming weekend and the next, and book the exam for when two consecutive mocks clear 80\%. Then do the more durable thing: keep the capstone repositories green. A platform you can still deploy from code in three months is worth more to your next employer than the badge.

\section{Milestone 6 Capstone: The Secure, Governed, Multi-Subscription Data Mesh}\label{sec:ch24-capstone}
\index{capstone project}\index{data mesh!capstone}\index{zero trust!capstone}

\begin{capstonebox}
\textbf{Mission:} Deploy the program's final architecture entirely from code: three subscriptions---ingestion with PII tokenization, analytics with Synapse and ADLS, governance with Purview---peered VNets, private endpoints everywhere, zero public access, dbt-gated loads, and a GitHub Actions pipeline that plans, tests, and applies under review.

\textbf{Estimated time:} 12--16 hours.

\textbf{What you will practice:}
\begin{itemize}
    \item Terraform modules per subscription with a locked remote state backend.
    \item Zero-trust networking: private endpoints, DNS, exfiltration protection (Chapter 21).
    \item Tokenized PII ingestion with event-driven processing (Chapters 13, 20, 23).
    \item Governance-as-default: Purview, dbt gates, budgets (Chapter 22).
\end{itemize}
\end{capstonebox}

\subsection*{Scenario}
The conglomerate from Chapter 22 sanctions the mesh pilot: Subscription A ingests streaming PII and tokenizes it with Functions and Key Vault; Subscription B hosts the Synapse and ADLS analytics estate; Subscription C runs Purview governance. All traffic is private; all deployment is code; all loads are quality-gated. You are the platform team.

\subsection*{Requirements}
\begin{itemize}
    \item \textbf{IaC:} one Terraform module per subscription; remote state with locking; nothing deployed by hand; \texttt{plan} on PRs, \texttt{apply} post-review.
    \item \textbf{Network:} VNet peering, private DNS zones, private endpoints on every data service; \texttt{public\_network\_access\_enabled = false} in code, verified from outside.
    \item \textbf{Ingestion:} Event Hubs $\rightarrow$ tokenization Function (Key Vault keys) $\rightarrow$ analytics estate; PHI never persists unmasked.
    \item \textbf{Quality:} dbt project with at least three models and five tests gating the tokenized load; poisoned-batch demonstration.
    \item \textbf{Governance:} sources registered and scanned in Purview; classifications evidenced before serving; cost budget per subscription with alerts.
    \item \textbf{Review:} the Chapter 22 four-page monthly review produced from this estate's real outputs.
\end{itemize}

\subsection*{Reference Architecture}
\begin{figure}[H]
    \centering
    \resizebox{\textwidth}{!}{%
    \begin{tikzpicture}[node distance=1.2cm]
        \node[stage] (suba) {Subscription A\\Ingestion\\Event Hubs + Functions\\(tokenize PII)};
        \node[stage, right=1.6cm of suba] (subb) {Subscription B\\Analytics\\Synapse + ADLS\\dbt-gated loads};
        \node[stage, below=1.4cm of suba] (subc) {Subscription C\\Governance\\Microsoft Purview};
        \node[doc, below=1.4cm of subb] (net) {Peered VNets\\private endpoints + DNS\\zero public access};
        \node[doc, right=1.6cm of subb] (cicd) {GitHub Actions\\plan $\rightarrow$ review $\rightarrow$ apply};
        \draw[arrow] (suba) -- node[above,font=\scriptsize]{VNet peering} (subb);
        \draw[arrow, dashed] (subc) -- (suba);
        \draw[arrow, dashed] (subc) -- (subb);
        \draw[arrow, dashed] (net) -- (subb);
        \draw[arrow, dashed] (net) -| (suba);
        \draw[arrow, dashed] (cicd) -- (subb);
    \end{tikzpicture}%
    }
    \caption{Final capstone: a three-subscription, zero-trust, code-deployed data mesh with governance as the default path.}
    \label{fig:ch24-capstone-arch}
\end{figure}

\subsection*{Implementation Steps}
\begin{enumerate}
    \item Write the three Terraform modules; configure the locked state backend; deploy Subscription C (governance) first.
    \item Peer the VNets; create private DNS zones; deploy endpoints with DNS zone groups---the Listing~\ref{lst:ch21-pep} pattern as code.
    \item Deploy the tokenization Function with Key Vault references; wire Event Hubs ingestion.
    \item Deploy Synapse with public access disabled; register and scan sources in Purview.
    \item Add the dbt project and the GitHub Actions gates; run the happy path and the poisoned batch.
    \item Prove zero public reachability from outside; produce the cost budgets and the monthly review.
\end{enumerate}

\subsection*{Deliverables}
\begin{itemize}
    \item The repository: Terraform modules, dbt project, workflows, and a README that deploys everything from clean.
    \item Evidence: plan/apply history, dbt test results (pass and blocked), the outside-fails network proof, Purview classifications.
    \item The cost report across three subscriptions and the four-page monthly review.
\end{itemize}

\subsection*{Acceptance Criteria}
\begin{itemize}
    \item \texttt{terraform destroy} followed by \texttt{apply} rebuilds the entire estate with no manual step.
    \item No resource is publicly reachable; all cross-subscription traffic flows over private endpoints.
    \item A poisoned tokenized batch is blocked by the dbt gate, quarantined, alerted, and audited.
    \item A destructive plan (e.g., deleting a storage account) is visible in PR review and blocked by process.
\end{itemize}

\subsection*{Stretch Goals}
\begin{itemize}
    \item Add a fourth domain subscription using only your modules and templates---the mesh handoff test.
    \item Implement one expand--contract migration on a served table with the lineage-gated contract step.
    \item Add Defender for Cloud and the Chapter 23 privacy calendar to the governance subscription.
\end{itemize}

\section{Key Takeaways for the Whole Program}
\begin{itemize}
    \item Storage decisions are architecture decisions; identity, networking, and lifecycle are designed, not defaulted.
    \item Serve logically when you can, provision when you must, and always measure bytes, seconds, and dollars.
    \item Pipelines are software: parameterized, tested, idempotent, and wired to fail loudly.
    \item Streams are state: watermarks, checkpoints, and exactly-once are engineering, not configuration.
    \item Models are data products: versioned features, gated promotion, and monitored drift.
    \item Governance, privacy, and cost are the operating layer that makes everything else survivable at scale.
    \item If it is not in code with a reviewed plan and a tested pipeline, it does not exist---that is the DataOps bar this book set out to teach.
\end{itemize}

\section{Exercises}
\begin{enumerate}
    \item \textbf{(Conceptual)} Why does Terraform state need locking, and what failure does the lease prevent?
    \item \textbf{(Conceptual)} Explain each expand--contract step and what breaks if step 4 precedes step 3.
    \item \textbf{(Conceptual)} Why run \texttt{terraform plan} on PRs but \texttt{apply} only after merge and review?
    \item \textbf{(Hands-on)} Complete the Thursday project with all three gate-proving pull requests.
    \item \textbf{(Hands-on)} Convert one hand-built lab environment to Terraform and verify zero drift on plan.
    \item \textbf{(Hands-on)} Add a dbt exposure documenting a Power BI dependency, then use lineage plus \texttt{ref()} to approve a contract step.
    \item \textbf{(Design)} Design the environment-promotion gates for the capstone estate: what evidence advances dev $\rightarrow$ test $\rightarrow$ prod.
    \item \textbf{(Design)} Write the ADR for Terraform versus Bicep for this estate.
    \item \textbf{(Design)} Build your final two-week certification plan from two mock-exam result sets.
    \item \textbf{(Interview)} Prepare a two-minute answer to ``walk me through how a column rename reaches production in your platform.''
\end{enumerate}
